\documentclass[9pt,landscape,a4paper]{article}
\usepackage[landscape,margin=0.6cm,top=0.6cm,bottom=0.6cm]{geometry}
\usepackage{fontspec}
\usepackage[brazilian]{babel}
\usepackage{xcolor}
\usepackage{multicol}
\usepackage{amsmath,amssymb}
\usepackage{booktabs}
\usepackage{tikz}
\usepackage{colortbl}

% --- TIPOGRAFIA ---
\setmainfont{Helvetica Neue}[
    Scale=0.75, 
    BoldFont={Helvetica Neue Bold},
    ItalicFont={Helvetica Neue Italic}
]
\newfontfamily\headerfont{Helvetica Neue Condensed Bold}[Scale=1.1]
\newfontfamily\titlefont{Helvetica Neue Condensed Bold}[Scale=1.5]

\pagestyle{empty}
\setlength{\parindent}{0pt}
\setlength{\columnsep}{0.6cm}

% --- CORES ---
\definecolor{csPrimary}{RGB}{0, 64, 128}   % Azul escuro elegante
\definecolor{csSecondary}{RGB}{0, 120, 180} % Azul médio
\definecolor{csAccent}{RGB}{200, 50, 50}    % Vermelho para alertas
\definecolor{rowA}{RGB}{255,255,255}
\definecolor{rowB}{RGB}{240,245,250}

% --- MACROS DE LAYOUT SEGURO (TIKZ PURO) ---
\newcommand{\csbox}[2]{%
\vspace{0.15cm}\noindent%
\begin{tikzpicture}
    \node[draw=csPrimary, fill=white, line width=0.8pt, rounded corners=3pt, inner sep=5pt, text width=\dimexpr\linewidth-10pt-1.6pt] {%
        {\color{csPrimary}\headerfont\large\MakeUppercase{#1}}\\[0.15cm]
        #2
    };
\end{tikzpicture}\par
}

\newcommand{\csalert}[2]{%
\vspace{0.15cm}\noindent%
\begin{tikzpicture}
    \node[draw=csAccent, fill=white, line width=1.5pt, rounded corners=3pt, inner sep=5pt, text width=\dimexpr\linewidth-10pt-3pt] {%
        {\color{csAccent}\headerfont\large\MakeUppercase{#1}}\\[0.15cm]
        #2
    };
\end{tikzpicture}\par
}

% --- MACROS DE LISTA COM ÍCONES ---
\newcommand{\cnum}[1]{%
    \raisebox{-0.5pt}{%
    \begin{tikzpicture}[baseline=(char.base)]
        \node[shape=circle,fill=csSecondary,text=white,inner sep=0.5pt,minimum size=3.5mm,font=\bfseries\scriptsize] (char) {#1};
    \end{tikzpicture}}%
}

\begin{document}

% HEADER DO CHEAT SHEET
\begin{tikzpicture}
    \node[draw=none, fill=csPrimary, rounded corners=4pt, inner sep=6pt, text width=\dimexpr\linewidth-12pt] {%
        {\textcolor{white}{\titlefont OXIMETRIA DE PULSO ($SpO_2$) - GUIA DE REFERÊNCIA RÁPIDA}}\\[0.05cm]
        {\textcolor{white!80}{\textbf{Autoria:} Nick Mark MD \quad |\quad \textbf{Tradução \& Adaptação:} Equipe ICU}}
    };
\end{tikzpicture}
\vspace{0.1cm}

\begin{multicols*}{3}

\csbox{Fisiologia Básica}{
A oximetria mede a saturação arterial de oxigênio de forma contínua usando dois princípios:

\vspace{0.15cm}
\textbf{1. Absorção Diferencial de Luz}\\
A sonda emite luz em dois comprimentos de onda:
\begin{itemize}
    \setlength{\itemsep}{1pt}
    \renewcommand{\labelitemi}{\textcolor{csAccent}{\tiny$\blacksquare$}}
    \item \textbf{Vermelho (660nm):} Absorvido pela \textbf{Hb Reduzida ($Hb$)}.
    \item \textbf{Infravermelho (940nm):} Absorvido pela \textbf{Hb Oxigenada ($HbO_2$)}.
\end{itemize}

\vspace{0.15cm}
\textbf{2. Fluxo Pulsátil (Pletismografia)}\\
Ignora a absorção constante (tecido, osso, sangue venoso). Calcula o oxigênio apenas no pulso arterial (fase sistólica).
}

\csbox{Curva de Dissociação da Hb}{
O $SpO_2$ (saturação) possui uma relação não-linear com o $PaO_2$ (oxigênio dissolvido no plasma).

\vspace{0.1cm}
\begin{itemize}
    \setlength{\itemsep}{1pt}
    \renewcommand{\labelitemi}{\cnum{1}}
    \item \textbf{Platô (Zona Segura):} Se $PaO_2 > 60$ mmHg, o $SpO_2 > 90\%$. Acima disso, o oxímetro é \textbf{insensível}. Um paciente com $PaO_2$ de 100 ou 500 mmHg lerá os mesmos 100\%.
    
    \renewcommand{\labelitemi}{\cnum{2}}
    \item \textbf{Porção Descendente (Perigo):} Quedas de $SpO_2$ abaixo de 90\% indicam redução DRÁSTICA do $PaO_2$ (<60 mmHg).
    
    \renewcommand{\labelitemi}{\cnum{3}}
    \item \textbf{Desvios da Curva:} 
    \begin{itemize}
        \renewcommand{\labelitemi}{\textcolor{csPrimary}{$\circ$}}
        \item \emph{Dir.:} Acidose, Febre (facilita entrega tecidual).
        \item \emph{Esq.:} Alcalose, Hipotermia.
    \end{itemize}
\end{itemize}
}

\csalert{ALERTA: Viés Racial e Hipoxemia}{
Oximetros modernos superestimam a saturação em pacientes com \textbf{pele escura}. 
\begin{itemize}
    \setlength{\itemsep}{1pt}
    \renewcommand{\labelitemi}{\textcolor{csAccent}{\bfseries !}}
    \item Risco 3x maior de \textbf{Hipoxemia Oculta} ($SpO_2$ parecendo normal na tela, mas $PaO_2$ crítico no sangue).
    \item \textbf{Conduta:} Correlacione sempre com a clínica. Diante de suspeita, colete Gasometria.
\end{itemize}
}

\csbox{Interferências Físicas}{
A precisão cai severamente se o $SpO_2$ for $< 80\%$ ou:

\textbf{Circulação Periférica Ruim}
\begin{itemize}
    \setlength{\itemsep}{1pt}
    \renewcommand{\labelitemi}{\textcolor{csSecondary}{\textbf{--}}}
    \item Choque, vasopressores em alta dose, ou manguito de PNI inflando no mesmo membro.
    \item \emph{Solução:} Troque para sensor central (orelha ou testa), que dão fluxo carotídeo e demoram menos para refletir quedas.
\end{itemize}

\textbf{Bloqueio Óptico e Ruído}
\begin{itemize}
    \setlength{\itemsep}{1pt}
    \renewcommand{\labelitemi}{\textcolor{csSecondary}{\textbf{--}}}
    \item \textbf{Esmalte:} Cor azul, preto e verde bloqueiam a luz.
    \item \textbf{Luz:} Refletores fortes saturam o fotodiodo. 
    \item \textbf{Movimento:} Tremores imitam o feixe pulsátil e confundem o aparelho.
\end{itemize}
}

\csbox{Interferência Tóxica (Hemoglobinas)}{
Certos agentes tóxicos alteram a hemoglobina, enganando as frequências de luz do sensor.

\rowcolors{1}{rowB}{rowA}
\begin{tabular}{p{2.2cm} p{1.6cm} p{3.2cm}}
\toprule
\textbf{Condição} & \textbf{Leitura} & \textbf{O que fazer?} \\
\midrule
\textbf{CO} (Monóxido) & Falso $\sim100\%$ mesmo c/ hipóxia. & Oxímetro é cego ao $CO$. Padrão: Gasometria c/ Co-Oximetria. \\
\textbf{Metemoglobina} & Fixa em $\sim85\%$. & Oximetria inútil. Co-Oximetria obrigatória. \\
\textbf{Azul Metileno} & Queda falsa. & Transitório (1-2 min). Aguardar. \\
\bottomrule
\end{tabular}
}

\csbox{Análise da Onda Pletismográfica}{
\textbf{Valide a qualidade da onda antes de tratar o número.}

\vspace{0.1cm}
\begin{center}
\begin{tikzpicture}[scale=0.9, font=\sffamily\tiny]
    \draw[gray, opacity=0.3] (0,0) grid (5,1.5);
    \draw[thick, csPrimary] (0,0.2) .. controls (0.2,0.2) and (0.3,1.3) .. (0.5,1.3) 
         .. controls (0.7,1.3) and (0.8,0.6) .. (0.9,0.7) 
         .. controls (1.0,0.8) and (1.4,0.2) .. (1.8,0.2)
         .. controls (2.0,0.2) and (2.1,1.3) .. (2.3,1.3)
         .. controls (2.5,1.3) and (2.6,0.6) .. (2.7,0.7) 
         .. controls (2.8,0.8) and (3.2,0.2) .. (3.6,0.2)
         .. controls (3.8,0.2) and (3.9,1.3) .. (4.1,1.3)
         .. controls (4.3,1.3) and (4.4,0.6) .. (4.5,0.7) 
         .. controls (4.6,0.8) and (5.0,0.2) .. (5.4,0.2);
         
    \node[anchor=south, text=csPrimary, font=\bfseries\scriptsize] at (2.3,1.3) {SÍSTOLE};
    \draw[->, csAccent, thick] (2.9, 1.1) -- (2.75, 0.75);
    \node[anchor=west, text=csAccent, font=\scriptsize] at (2.9,1.1) {Incisura Dicrótica};
\end{tikzpicture}
\end{center}

\vspace{0.1cm}
\begin{itemize}
    \setlength{\itemsep}{2pt}
    \renewcommand{\labelitemi}{\cnum{1}}
    \item \textbf{Pico Sistólico:} Deve acompanhar o batimento real.
    \renewcommand{\labelitemi}{\cnum{2}}
    \item \textbf{Incisura Dicrótica:} lombada na descida representa fechamento da valva aórtica. Atesta boa perfusão arterial.
    \renewcommand{\labelitemi}{\cnum{3}}
    \item \textbf{Validação FC:} A Frequência Cardíaca ($bpm$) do oxímetro DEVE bater com o ECG. Do contrário, o sinal é lixo.
\end{itemize}
}

\csalert{O Padrão-Ouro}{
Na dúvida (ex: choque, hipoxemia), \textbf{não confie no monitor}.\\[0.1cm]
Sempre colete a \textbf{Gasometria Arterial} garantindo Co-Oximetria em suspeitas tóxicas.
}

\end{multicols*}

\end{document}
